\clearpage
\VerdierPhysicalPageAnchor{0085}%
\VerdierSemanticAnchor{ch1-3-1-4-homotopy-definition-continuation}
を $\operatorname{tr}_1(C)$ の二つの射とする。$(g,h,a)$ と
$(g',h',a')$ を結ぶ\emph{ホモトピー}とは三つ組 $(b,b',\alpha)$ であり、
$b\colon g\rightsquigarrow g'$ と $b'\colon h\rightsquigarrow h'$ は
ホモトピーで、

\VerdierSemanticAnchor{ch1-3-1-4-second-order-condition}
\begin{equation*}
  \alpha\colon f'\star b+a'\rightsquigarrow b'\star f+a
\end{equation*}
は二階ホモトピー (2.5.12) である。

\VerdierSemanticAnchor{ch1-3-1-4-homotopic-morphisms}
$\operatorname{tr}_1(C)$ の二つの射は、それらを結ぶホモトピーが存在するとき、
\emph{ホモトピック}であるという。
\VerdierSemanticAnchor{ch1-3-1-4-equivalence-relation}
「$(g,h,a)$ と $(g',h',a')$ はホモトピックである」という関係は、射の直和と
合成に両立する同値関係である。
\VerdierSemanticAnchor{ch1-3-1-4-tr1-homotopy-category}
したがって、対象を $\operatorname{tr}_1(C)$ の対象、射を
$\operatorname{tr}_1(C)$ の射のホモトピー類とする圏
$\operatorname{Tr}_1(C)$ を導入できる。

\VerdierSemanticAnchor{ch1-3-1-5}
\noindent\textbf{命題 3.1.5.}\quad
\emph{
$D,D'\colon (X\xrightarrow{f}Y)\rightrightarrows
(X'\xrightarrow{f'}Y')$ を $\operatorname{tr}_1(C)$ の二つの射とする。
$D$ と $D'$ が $\operatorname{tr}_1(C)$ でホモトピックならば、$c(D)$ と
$c(D')$ は $\operatorname{comp}(C)$ でホモトピックである。}

\VerdierSemanticAnchor{ch1-3-1-5-proof-setup}
$D=(g,h,a)$、$D'=(g',h',a')$ とする。$D$ と $D'$ が
$\operatorname{tr}_1(C)$ でホモトピックであると仮定する。このとき、次を満たす
ホモトピー $b,b'$ が存在する：

\VerdierSemanticAnchor{ch1-3-1-5-eq1}
\begin{equation}
  \begin{aligned}
    g'-g&=bd+db,\\
    h'-h&=b'd+db'.
  \end{aligned}
  \tag{3.1.5.1}
\end{equation}

（記号を軽くするため、この計算では各成分の次数を表す上付き添字を省いた。）
\VerdierSemanticAnchor{ch1-3-1-5-second-order-homotopy}
さらに、次を満たす二階ホモトピー $\alpha$ が存在する：
\begin{equation*}
  \alpha d-d\alpha=b'f+a-(f'b+a').
\end{equation*}

\VerdierSemanticAnchor{ch1-3-1-5-cone-homotopy}
直ちに計算すると
\begin{equation*}
  \begin{pmatrix}
    b' & \alpha\\
    0  & -b
  \end{pmatrix}
\end{equation*}
が $c(D)$ と $c(D')$ を結ぶホモトピーであることがわかる。

\VerdierSemanticAnchor{ch1-3-1-6}
\noindent\textbf{系 3.1.6.}\quad
\emph{関手}
\begin{equation*}
  c\colon\operatorname{tr}_1(C)\longrightarrow\operatorname{comp}(C)
  \qquad\text{および}\qquad
  s\colon\operatorname{comp}(C)\longrightarrow\operatorname{tr}_1(C)
\end{equation*}
\emph{は $\operatorname{tr}_1(C)$ と $\operatorname{comp}(C)$ のホモトピーに
両立し、商をとることにより、やはり}

\VerdierSemanticAnchor{ch1-3-1-6-descended-functors}
\begin{equation*}
  c\colon\operatorname{Tr}_1(C)\longrightarrow K(C)
  \qquad\text{および}\qquad
  s\colon K(C)\longrightarrow\operatorname{Tr}_1(C)
\end{equation*}
\emph{と書かれる関手を定める。}
\VerdierSemanticAnchor{ch1-3-1-6-adjunction}
\emph{関手 $c\colon\operatorname{Tr}_1(C)\longrightarrow K(C)$ は関手
$s\colon K(C)\longrightarrow\operatorname{Tr}_1(C)$ の左随伴である。}

\VerdierSemanticAnchor{ch1-3-1-6-proof}
これは (3.1.5) と (3.1.3) から形式的に従う。

\VerdierSemanticAnchor{ch1-3-1-7}
\noindent\textbf{3.1.7.}\quad
\VerdierSemanticAnchor{ch1-3-1-7-s-prime}
$s'\colon\operatorname{comp}(C)\longrightarrow\operatorname{tr}_1(C)$ を
関手 $X\longmapsto(X\longrightarrow 0)$ とする。
\VerdierSemanticAnchor{ch1-3-1-7-comp-translation}
関手 $cs'\colon\operatorname{comp}(C)\longrightarrow\operatorname{comp}(C)$
は $\operatorname{comp}(C)$ のシフト関手に等しい。
\VerdierSemanticAnchor{ch1-3-1-7-descended-s-prime}
関手 $s'$ は $\operatorname{comp}(C)$ のホモトピックな二つの射を
$\operatorname{tr}_1(C)$ のホモトピックな射に送り、したがって、やはり $s'$ と
書かれる関手 $s'\colon K(C)\longrightarrow\operatorname{Tr}_1(C)$ を定める。
\VerdierSemanticAnchor{ch1-3-1-7-k-translation}
関手 $cs'\colon K(C)\longrightarrow K(C)$ は $K(C)$ のシフト関手に等しい。
